\documentclass{beamer}
\usetheme{boxes}
\usecolortheme{default}
\usepackage[T1]{fontenc}
\usepackage{amsmath,amssymb}
\usepackage{csquotes}
\usepackage{array}
\usepackage{colortbl}
\setbeamertemplate{navigation symbols}{}


% CUSTOM FOOTER CONFIGURATION (Page number / Total pages)
\setbeamertemplate{footline}{%
  \leavevmode%
  \hbox{%
    \begin{beamercolorbox}[wd=\paperwidth,ht=2.25ex,dp=1ex,right,rightskip=2ex]{footline}%
      \usebeamerfont{page number in head/foot}%
      \insertframenumber{} / \inserttotalframenumber{}
    \end{beamercolorbox}%
  }%
  \vskip0pt%
}


% Style: C:/mytex/2026-Krakow-pres.tex.
% Source: C:/mytex/2026-trade.tex, dated July 29, 2026.
% Focus on the paper's computable countermarket construction.
% Discounted price: slide 5. Explicit worked example: slide 9.
% Private randomization: slides 12--13. Permutation stress test: slide 20.
% Former financial digressions remain removed.
% Exactly 21 frames. All numerical paths are illustrative, not market data.
% Self-contained; compile twice with pdflatex 2026-ECONVN-pres.tex.
\newcommand{\doiref}[2]{%
  \href{https://doi.org/#1}{\textcolor{blue!65!black}{#2}}}
\newcommand{\sourcefoot}[1]{%
  \par\vfill{\fontsize{6.3}{7.5}\selectfont\color{black!65}Sources: #1\par}}

\title{No Trading Strategy Can Win\\on Every Price Path}
\subtitle{A Computable Countermarket for Every Trader}
\author{\texorpdfstring{Karl Svozil\\ITP, TU Wien, Vienna, Austria}%
{Karl Svozil, ITP, TU Wien, Vienna, Austria}}
\date{ECONVN 2026/2027, HUB/AJEB,\\ Ho Chi Minh City, Vietnam\\
12-14 October 2026\\
{\footnotesize \url{https://svozil.github.io/publications/2026-ECONVN-pres.pdf} }
\\
based on
\\
\url{https://doi.org/10.48550/arXiv.2604.13334}
}

\begin{document}

% Slide 1.
\begin{frame}
\titlepage
\end{frame}

% Slide 2.
\begin{frame}{Finite profits do not establish universal profitability}
\small
\begin{block}{A trader can make money over a finite period}
Accurate information, luck, or a continuing trend can produce profits.
Exuberance, innovation, and improvements in production can support a rise.
\end{block}
\begin{exampleblock}{The same kind of bet can also lose}
Information may be false. Exuberance may turn into panic.
Expected gains from innovation or cost cutting may fail to reach shareholders.
\end{exampleblock}
\begin{alertblock}{The claim that fails}
\enquote{One trading rule that earns strict profit on every possible
price history.}
\par\medskip
The reason is constructive (algorithmically provable): use the proposed rule to build its
\emph{countermarket}, a price record on which it cannot win.
\end{alertblock}
\end{frame}


% Slide 3.
\begin{frame}{The trader chooses from the history so far}
\small
\begin{block}{A trading rule}
The rule observes past and current prices and chooses a position before
the next price arrives. It cannot use a future price that it has not observed.
\end{block}
\begin{exampleblock}{Three choices}
\begin{description}
\setlength{\itemsep}{0.55em}
\item[Long:] own shares and gain when their price rises.
\item[Short:] borrow and sell shares, intending to buy them back later;
gain when their price falls.
\item[Wait:] hold neither a long nor a short position; earn zero trading gain.
\end{description}
\end{exampleblock}
\begin{alertblock}{The mathematical setting}
The rule is an algorithm that returns an action on every finite history
(\emph{total and computable}), without random choices (\emph{deterministic}).
Position sizes are ratios of integers (\emph{rational} amounts).
\end{alertblock}
\end{frame}

\begin{frame}
\hskip 4 cm
\begin{center}
\Large\color{blue}
Countermarket construction
\end{center}
\end{frame}

% Slide 4.
\begin{frame}[shrink=8]{Cantor's diagonal idea}
\small
Diagonal reasoning also appears in the work of G\"odel and Turing.
Suppose there exists a \emph{complete} list containing \emph{all} infinite sequences of zeros and
ones. Its beginning might look like this:
\[
\begin{array}{c|ccccc}
 & \text{bit 1}&\text{bit 2}&\text{bit 3}&\text{bit 4}&\cdots\\\hline
\text{row 1}&\textcolor{red}{0}&1&1&0&\cdots\\
\text{row 2}&1&\textcolor{red}{1}&0&1&\cdots\\
\text{row 3}&0&1&\textcolor{red}{0}&0&\cdots\\
\text{row 4}&1&0&1&\textcolor{red}{1}&\cdots\\
\vdots&\vdots&\vdots&\vdots&\vdots&\ddots
\end{array}
\]
\begin{exampleblock}{Change the diagonal entries}
Read the highlighted bits and change each $0$ to $1$, each $1$ to $0$.
The new sequence begins $1,0,1,0,\ldots$.
\end{exampleblock}
\begin{alertblock}{Why the list cannot be complete}
The new sequence differs from row 1 at bit 1, from row 2 at bit 2,
and from every row at its diagonal position. It is absent from the list.
\end{alertblock}
\sourcefoot{Metamathematics; eg, Noson S. Yanofsky,
\doiref{10.2178/bsl/1058448677}{\emph{A Universal Approach to Self-Referential Paradoxes, Incompleteness and Fixed Points} (2003)}.}
\end{frame}

% Slide 5.
\begin{frame}[shrink=5]{Discounted price: value compared with the bank}
\small
\begin{block}{Definition}
A \emph{discounted price} is the asset's price divided by the current
value of one unit of the comparison bank account.
It tells us how many such bank units the asset is worth.
The model treats this account as risk-free.  All money amounts below are in a fixed currency; eg, Vietnamese dong (VND).
\end{block}
\begin{exampleblock}{Example: one bank unit is a deposit of 100 VND}
Suppose the account earns $10\%$ over the year, whereas the share sells at 121 VND.
\par\medskip
\centering
\renewcommand{\arraystretch}{1.5}
\begin{tabular}{@{}lrrc@{}}
 & Bank unit & One share & Discounted price\\\hline
At the start & $100$ & $100$ & $100/100=1$\\
After one year & $110$ & $121$ & $121/110=1.10$
\end{tabular}
\end{exampleblock}
\begin{alertblock}{What the number means}
Selling the share now buys $1.10$ bank units instead of $1$.
The share gained $10\%$ \emph{relative to the bank account}.
Had the share risen only to $110$, its discounted price would still be $1$.
(This adjusts for bank interest, not inflation.)
\end{alertblock}
\end{frame}

% Requires \usepackage{colortbl} in the preamble.

% Slide 6.
\begin{frame}{From Cantor's idea to a countermarket}
\small

\begin{tabular}{@{}
  >{\raggedright\arraybackslash}p{0.485\textwidth}
  @{\hspace{0.03\textwidth}}
  >{\raggedright\arraybackslash}p{0.485\textwidth}
  @{}}
\textcolor{blue}{Cantor's construction}
&
\textcolor{green!50!black}{The trading construction}
\\[0.8em]

\rowcolor{black!15}
Start with a purportedly complete enumeration.
&
Start with a proposed universally profitable rule.
\\[1em]


Use that list itself to construct something that it misses.
&
Use that rule itself to construct a price path on which it fails.
\\[1em]

\rowcolor{black!15}
Changing the selected bit defeats the completeness claim.
&
Opposing each position defeats the universal-profit claim.
\end{tabular}

\medskip
\begin{alertblock}{The common idea is diagonal opposition}
The counterexample depends on the proposed universal object.
Here we construct a market for one fixed trader; we do not need to list
all trading programs.
\end{alertblock}
\end{frame}

% Slide 7.
\begin{frame}{A Computable Countermarket for Every Trader}
\small
\begin{block}{The central result developed in the paper}
For every deterministic program that returns a rational share position in
finite time on every finite price history, there is a fixed computable price
path whose prices remain strictly above zero and on which accumulated
discounted trading gain is never positive at any finite date.
\end{block}
\begin{exampleblock}{The stronger statement about active trading}
Every nonzero position loses on the constructed path.
After any period of active trading, its accumulated gain is strictly
negative. A trader that always waits earns zero.
\end{exampleblock}
\begin{alertblock}{What we are measuring}
Prices are discounted as on slide 5: gain is measured relative to the bank
account. Trading is \emph{self-financing}: new positions use the shares and
bank balance already held, with no outside cash added. The simple proof
omits trading fees, which would further reduce gains.
\end{alertblock}
\sourcefoot{KS, \doiref{10.48550/arXiv.2604.13334}{\emph{No Trading Strategy Can Win on Every Price Path},
subsection \enquote{A Computable Countermarket for Every Trader} (2026)}.}
\end{frame}

% Slide 8.
\begin{frame}[shrink=5]{The countermarket: a constructive proof}
\small
\begin{block}{Use the proposed trader's own program}
Calculate its position from the history constructed so far. Then choose
the next price move according to this table:
\end{block}
\begin{exampleblock}{Opposition at every trading date}
\centering
\renewcommand{\arraystretch}{1.5}
\begin{tabular}{lll}
Trader's choice & Countermarket's move & Trading result\\\hline
\rowcolor{black!15}
Hold or buy shares & Price falls & Loss\\
Sell borrowed shares & Price rises & Loss\\
\rowcolor{black!15}
Hold no position & Price stays the same & Zero
\end{tabular}
\end{exampleblock}
\begin{alertblock}{Why this proves the claim}
Choose a fixed, small percentage rise or fall that keeps prices positive,
and repeat.
Adding losses and zeros can never produce a positive gain.
For each fixed trader this gives one fixed computable price path, whereby the countermarket is constructed from this trader's program.
Any finite part of its price history can be calculated before trading begins.
It need not defeat other traders.
\end{alertblock}
\end{frame}

% Slide 9.
\begin{frame}[shrink=5]{Worked example: we construct the adverse prices}
\small
\begin{block}{The trader's rule is fixed first}
The rule is: buy for one period, sell short for the next, then wait.
All prices are discounted prices; gains and losses are measured in
bank-account units.
\end{block}
\begin{exampleblock}{Our construction, step by step}
\begin{enumerate}
\setlength{\itemsep}{0.55em}
\item The trader buys one share for $100$.\newline
We choose $90$ as the next price. Selling the share then loses $10$.
\item The trader borrows a share and sells it for $90$.\newline
We choose $99$ as the next price. Buying it back to return it loses $9$.
\item The trader holds no position.\newline
We leave the price at $99$. Its trading gain is zero.
\end{enumerate}
\end{exampleblock}
\begin{alertblock}{Why this is a counterexample}
This constructed history makes the example rule lose.
It therefore refutes any claim that this rule earns a strictly positive
gain on every price history.
\end{alertblock}
\end{frame}

% Slide 10.
\begin{frame}{Why the countermarket is a fixed computable path}
\small
\begin{block}{Put the trader's code inside a price-generating program}
\begin{enumerate}
\setlength{\itemsep}{0.45em}
\item Start with a positive price and the history built so far.
\item Run the trader's program on that history to obtain its position.
\item Choose a small price fall for a long position, a small rise for a
short position, and no change for zero holdings.
\item Append that price to the history and repeat.
\end{enumerate}
\end{block}
\begin{exampleblock}{Each requested finite part can be calculated}
The trader returns its position in finite time. Rational positions, expressed
as ratios of integers, can be compared with zero exactly. Each next price is therefore
computable.
\end{exampleblock}
\begin{alertblock}{The market does not have to observe live orders}
For a fixed trader, the generator is fixed. Any finite portion can be
computed before trading begins, although computation may be slow.
\end{alertblock}
\end{frame}

% Slide 11.
\begin{frame}{Every trader has its own countermarket}
\small
\begin{block}{The order of the claim matters}
Proposed guarantee: \enquote{There is one trader that profits on every path.}
\par\medskip
Countermarket result: \enquote{For each proposed trader, construct a path
on which every nonzero position loses.}
\end{block}
\begin{exampleblock}{Another trader may win on that same record}
For one price rise, a trader holding one share gains and a trader shorting
one share loses. A price fall reverses their fortunes.
The construction does not claim to make every other trader lose on that path.
\end{exampleblock}
\begin{alertblock}{The stronger contribution}
A constant discounted price already gives every trader zero trading gain.
Our construction makes every active position of the selected trader lose.
The resulting path need not be typical or an economic equilibrium.
The following results address separate questions.
\end{alertblock}
\end{frame}

% Slide 12.
\begin{frame}[shrink=8]{Randomized trading: what does the trader randomize?}
\small
\begin{block}{A private coin chooses the position}
After observing the past, the trader may toss a private coin:
heads means own one share, tails means sell one borrowed share.
The coin chooses the \emph{position}.
The subsequent price movement determines whether that position gains or loses.
\end{block}

\begin{alertblock}{Why the earlier diagonal construction needs qualification}
For a deterministic trader, its program and the price history determine
its next position.
With private randomness, we also need the outcome of its coin toss.
The preceding countermarket construction therefore cannot calculate and
oppose each hidden choice from the program and price history alone.
\end{alertblock}

\begin{exampleblock}{A separate argument using an independent market}
We instead let the market choose its price movements independently of
the trader's orders. Such a market is called \emph{passive}: it does not
react to those orders.
The trader may still randomize, so both sides can randomize simultaneously.
\par\smallskip
We now calculate \emph{expected gain}: the probability-weighted average
of the possible gains and losses.
\end{exampleblock}
\sourcefoot{KS, subsection \enquote{Worst-Case Value and Private Randomization}.}
\end{frame}

% Slide 13.
\begin{frame}[shrink=8]{An independent fair market gives zero expected gain}
\small
\begin{block}{The market's coin chooses the price movement}
The trader first chooses its position, possibly at random.
An independent fair coin then moves the discounted price from $100$
to $90$ or $110$, each with probability one half.
The market's coin is independent of the past and of the trader's private coin.
\end{block}

\begin{exampleblock}{Average over the market's two equally likely moves}
\centering
\renewcommand{\arraystretch}{1.3}
\setlength{\tabcolsep}{8pt}
\begin{tabular}{@{}lrrc@{}}
Position & Price $90$ & Price $110$ & Expected gain\\\hline
\rowcolor{black!15}
Own one share & $-10$ & $+10$ & $0$\\
Short one share & $+10$ & $-10$ & $0$\\
\rowcolor{black!15}
Hold no position & $0$ & $0$ & $0$
\end{tabular}
\end{exampleblock}

\begin{alertblock}{The result applies to deterministic and randomized traders}
Every chosen position has expected gain zero, however the trader selected it.
Repeat with independent fair choices of a $10\%$ rise or fall from each
current price. Expected total gain over any fixed finite number of trades
is then zero, provided gains and losses have finite averages.
\par\smallskip
Individual runs may win or lose. Randomizing the position cannot guarantee
positive expected gain in \emph{every} market.
\end{alertblock}
\sourcefoot{KS, proposition \enquote{Private Randomization Does Not Create a Guarantee}.}
\end{frame}

\begin{frame}
\hskip 4 cm
\begin{center}
\Large\color{blue}Some related ideas
\end{center}
\end{frame}


% Slide 14.
\begin{frame}{Gold: can history reveal the market's rule?}
\small
\begin{block}{Learning in the limit}
After each new observation, a learner proposes a program explaining the
whole sequence. Success means eventually settling on one correct program
and never changing it again. The learner need not know when it has succeeded.
\end{block}
\begin{alertblock}{Gold's negative result}
No algorithm achieves this for \emph{every} computable sequence of zeros
and ones. For each learner, some computable history defeats this requirement.
Failure means no eventual correct identification, rather than a specified
number of losing trades.
\end{alertblock}
\begin{exampleblock}{Restrictions can make learning possible}
If an algorithm can list the allowed rules, and each rule computes every
requested term in finite time, identification in the limit is possible.
What matters is which market rules the learner admits.
\end{exampleblock}
\sourcefoot{E.\,M. Gold,
\doiref{10.2307/2270580}{\emph{Limiting recursion} (1965)};
\doiref{10.1016/S0019-9958(67)91165-5}{\emph{Language identification in the limit} (1967)}.
Related prediction construction: \doiref{10.48550/arXiv.cs/0606070}{Legg, Is there an Elegant Universal Theory of Prediction? (2006)}.}
\end{frame}

% Slide 15.
\begin{frame}{Turing: will the price ever rise?}
\small

\begin{block}{An artificial market controlled by a program}
Run any program with fixed input, one instruction per day. Start the price at $1$.
\begin{itemize}
\setlength{\itemsep}{0.2em}
\item While the program keeps running, the price stays at $1$.
\item If the program stops, the price becomes $2$ and stays there.
\end{itemize}
\end{block}

\begin{exampleblock}{The price on a specified day is computable}
\centering
\renewcommand{\arraystretch}{1.2}
\begin{tabular}{@{}lccccc@{}}
Program's behaviour & Day 1 & Day 2 & Day 3 & Day 4 & $\cdots$\\\hline
Stops at instruction 3 & $1$ & $1$ & $2$ & $2$ & $\cdots$\\
Never stops            & $1$ & $1$ & $1$ & $1$ & $\cdots$
\end{tabular}
\par\smallskip
\raggedright
For day $100$, simulate at most $100$ instructions.
If the program has stopped, report $2$; otherwise, report $1$.
This calculation always finishes.
\end{exampleblock}

\begin{alertblock}{But will the price \emph{ever} rise?}
The price ever rises exactly when the program eventually stops.
This is the \emph{Halting Problem}: no algorithm can always finish with the
correct yes-or-no answer for every program and input, even with its source code.
\end{alertblock}

\sourcefoot{A.\,M. Turing,
\doiref{10.1112/plms/s2-42.1.230}{\emph{On computable numbers} (1936--37)}.
Market construction: the paper.}
\end{frame}

% Slide 16.
\begin{frame}{Rad\'o's Busy Beaver: how long must we wait?}
\small
\begin{block}{A finite maximum with no computable general bound}
Fix a \emph{universal} programming language, capable of simulating any
ordinary program. Include any required input in its program description.
Among programs of at most $n$ bits that eventually stop, take the longest
running time. This is the runtime
\emph{Busy Beaver} value, written $\mathrm{BB}(n)$.
Here $n$ measures program length, not elapsed time.
\end{block}
\begin{exampleblock}{Return to the price $1$ that may become $2$}
A short program can run for an extraordinarily long time before stopping.
Its market can therefore remain flat through a long observation period and
change only later.
\end{exampleblock}
\begin{alertblock}{Why no waiting rule suffices}
No computable deadline based only on program length covers every eventual
halt. Such a deadline would solve Turing's problem: wait until it expires.
Some finite cases can be settled, but they give no general extrapolation rule.
\end{alertblock}
\sourcefoot{T. Rad\'o,
\doiref{10.1002/j.1538-7305.1962.tb00480.x}{\emph{On non-computable functions} (1962)};
G.\,J. Chaitin,
\doiref{10.1007/978-1-4612-4808-8_28}{\emph{Computing the busy beaver function} (1987)}.}
\end{frame}

% Slide 17.
\begin{frame}{Martin-L\"of: randomness still permits finite wins}
\small
\begin{block}{Random relative to which model?}
Use fair-coin outcomes, written as zeros and ones.
A test might flag an initial run of only zeros: its probability is $1/2$
for one zero, $1/4$ for two, and so on.
Informally, a \emph{Martin-L\"of-random} infinite record cannot remain
rejected at every stage of any algorithmic test with guaranteed
probability bounds shrinking to zero.
\end{block}
\begin{alertblock}{The betting consequence}
The trader may never risk more than its available capital, so its balance
stays nonnegative on every possible history. In this fair betting game,
no computable strategy can make its capital unbounded on such a record.
\emph{Fair} means that, conditional on the history, the average next capital
equals the current capital.
\end{alertblock}
\begin{exampleblock}{What remains possible}
Finite winning streaks, apparent trends, and correct forecasts still occur.
Every finite record has both computable and algorithmically random
continuations. A finite sample alone cannot certify randomness.
\end{exampleblock}
\sourcefoot{P. Martin-L\"of,
\doiref{10.1016/S0019-9958(66)80018-9}{\emph{The definition of random sequences} (1966)};
J. Ville, \href{https://www.numdam.org/item/THESE_1939__218__1_0/}{\emph{\'Etude critique de la notion de collectif} (1939)}.}
\end{frame}

% Slide 18.
\begin{frame}{Ramsey: order appears even without a predictive law}
\small
\begin{block}{An unavoidable pattern}
Join every pair among six objects and color each link red or blue.
Whatever the coloring, three objects have all three links the same color.
This is a basic finite case of Ramsey's theorem.
\end{block}
\begin{exampleblock}{Repeated symbols at equally spaced dates}
Van der Waerden's theorem gives a related result: in a sufficiently long
record with finitely many symbols, some symbol repeats at equally spaced
positions for any prescribed finite number of repetitions.
\end{exampleblock}
\begin{alertblock}{The finite-pattern trap}
The selected dates need not be consecutive, the returns need not be positive,
and the pattern need not continue. Calude and Longo emphasize that large
data sets force some regularities even when the records resist compression.
Such patterns coexist with Martin-L\"of randomness.
\end{alertblock}
\sourcefoot{Graham, Rothschild \& Spencer, \emph{Ramsey Theory} (1990);
\doiref{10.1038/scientificamerican0790-112}{Graham \& Spencer (1990)};
\doiref{10.1007/s10699-016-9489-4}{Calude \& Longo, \emph{The deluge of spurious correlations in big data} (2017)}.}
\end{frame}

% Slide 19.
\begin{frame}{Time reversal as a search for counterexamples}
\small
\begin{block}{Hold the trading rule fixed, change the record}
Test it on the illustrative prices $100,\ 110,\ 120$ and then rerun it
from scratch on $120,\ 110,\ 100$.
At every date it may use only the history available on that run.
Reversing the original orders would give it information from the future.
\end{block}
\begin{alertblock}{One failure refutes an all-path claim}
If both records belong to the claimed market class, a nonpositive gain on
either refutes strict profit on every path.
Reversal alone does not prove that a rule must fail on one of the pair.
It supplies a stress test, whereas diagonalization supplies the counterexample.
\end{alertblock}
\sourcefoot{KS, \enquote{A Practical Robustness Diagnostic: Time Reversal},
Reversal Falsification Test.}
\end{frame}

% Slide 20.
\begin{frame}[shrink=8]{Beyond time reversal: reshuffling the record}
\small
\begin{block}{Algorithmic and random permutations}
A \emph{permutation} rearranges the same finite record of discounted price
levels. Use a fixed algorithm or a random shuffle. Rerun the trader from
scratch, giving it only the past available in the new order.
\end{block}
\begin{alertblock}{A necessary condition for a universal profit guarantee}
A strategy claiming strict profit on \emph{every} permitted path
\textbf{must remain profitable under every reshuffling that stays within
its claimed market class}. The amount of profit may change.
One zero or negative gain refutes that universal claim.
\end{alertblock}
\begin{exampleblock}{What this requirement means}
A strategy exploiting a particular trend may fail when reshuffling destroys
that trend. Such failure limits its scope. Conversely, passing a finite
collection of shuffles cannot establish universal profitability.
\end{exampleblock}
\medskip
\emph{Literary analogy:} Greg Egan's \emph{Permutation City} (1994)
explores the ordering of computational states.
\sourcefoot{Greg Egan, \href{https://www.gregegan.net/PERMUTATION/Permutation.html}{\emph{Permutation City} (1994)};
\href{https://www.gregegan.net/PERMUTATION/FAQ/FAQ.html}{Egan's \emph{Dust Theory FAQ}, questions 1--2}.}
\end{frame}

% Slide 21.
\begin{frame}{The contribution: an explicit countermarket}
\small
\begin{block}{Computable diagonal countermarket}
For every total deterministic computable trading rule with rational positions,
there exists a positive computable price path on which its self-financing
gain is nonpositive at every finite horizon.
\end{block}
\begin{alertblock}{Scope of the result}
The countermarket depends on the chosen rule. Every nonzero position incurs
a loss, while zero holdings yield zero gain. Thus no such rule guarantees
strictly positive discounted gain on every positive computable price path.
\end{alertblock}
\begin{exampleblock}{Finite profitability}
Information, luck, or favorable market conditions can produce finite profits.
These remain compatible with adverse continuations and do not establish
a universal guarantee.
\end{exampleblock}
\vfill
\end{frame}


\begin{frame}
\hskip 4 cm
\begin{center}
\Large\color{blue}Thank you for your attention
\end{center}
\end{frame}

\end{document}
